% ===================================================================
% FST-IV: Ein skaleninvariantes Stabilitätsprinzip von Teilchen bis zur Kosmologie
%         -- Ein Überblick über das Functional Stability Theory Programm
% Target journal: Foundations of Physics / Entropy
% Status: v1 (programmatischer Überblick, struktureller und reduktiver Beitrag)
% Author: Lukas Geiger, Bernau, Germany
% Versionsgeschichte:
%   v0 -- 2026-04 (initiales Skeleton, struktureller Entwurf)
%   v1 -- 2026-05-12 (Series-Hub: Komponententabelle, Building-Blocks-Status,
%                     Bibitems mit Zenodo-DOIs, AI-Disclosure L3,
%                     EN/DE-Parität)
%   v2 -- 2026-05-14 (Quellencheck: veraltete Zenodo-Record-DOIs durch
%                     Concept-DOIs ersetzt; externe Referenzen mit
%                     verifizierten DOI-Resolvern ergänzt)
% ===================================================================

\documentclass[11pt,a4paper]{article}

% --- Packages ------------------------------------------------------
\usepackage[utf8]{inputenc}
\usepackage[T1]{fontenc}
\usepackage{lmodern}
\usepackage[ngerman]{babel}
\usepackage{amsmath,amssymb,amsthm}
\usepackage{graphicx}
\usepackage{xurl}
\input{glyphtounicode}
\pdfgentounicode=1
\usepackage[unicode=true,pdfdisplaydoctitle=true]{hyperref}
\usepackage{booktabs}
\usepackage{array}
\usepackage{tabularx}
\usepackage[margin=2.4cm]{geometry}
\usepackage{enumitem}
\usepackage{xcolor}
\newcolumntype{Y}{>{\raggedright\arraybackslash}X}
\setlength{\emergencystretch}{3em}
\hfuzz=5pt
\makeatletter
\renewcommand*{\HyperDestNameFilter}[1]{de.#1}
\makeatother
\hypersetup{
  pdftitle={Functional Stability Theory - Ein programmatischer Hub: Universelle Konvexitäts-Eindeutigkeit über Skalen hinweg},
  pdfauthor={Lukas Geiger},
  pdfsubject={Programmatischer Überblick zur Functional Stability Theory},
  pdfkeywords={functional stability theory, universelle konvexitäts-eindeutigkeit, renormierungsgruppe, thermodynamik, kosmologie}
}

% --- Theorem environments ------------------------------------------
\newtheorem{axiom}{Axiom}
\newtheorem{definition}{Definition}
\newtheorem{lemma}{Lemma}
\newtheorem{proposition}{Proposition}
\newtheorem{conjecture}{Vermutung}
\theoremstyle{remark}
\newtheorem*{remark}{Bemerkung}

% --- Title ---------------------------------------------------------
\title{Functional Stability Theory \,---\,\\
Ein programmatischer Hub:\\
Universelle Konvexitäts-Eindeutigkeit\\
über Skalen hinweg\\[0.5em]
\large\textit{Von Teilchen bis zur Kosmologie}}

\author{Lukas Geiger\\
\textit{Unabhängiger Forscher, Bernau, Deutschland}\\
ORCID: \url{https://orcid.org/0009-0005-7296-1534}}

\date{Mai 2026}

\begin{document}
\maketitle

\begin{abstract}
\noindent
Functional Stability Theory (FST) ist ein Forschungsprogramm, das ein einziges Stabilitätsprinzip -- das \emph{Universal Convexity-Uniqueness Kriterium} (UCU) -- als vereinheitlichenden Rahmen für emergente Strukturen über physikalische Skalen hinweg vorschlägt: von thermodynamischen Gleichgewichten fundamentaler Parameter (FST-I) über autokatalytische Chemie (FST-II) und biologische evolutionäre Stabilität (FST-III) bis hin zu kosmologischer Dunkler Energie und Turbulenz. Dieser Überblick synthetisiert das Programm. Das universelle Spiel $(x_L, G_L, D_L)$ wird zusammen mit dem Prinzip der maximalen Entropieproduktion (MEPP) als Anfangsbedingungsaxiom skizziert, und das Drei-Lemma-Endspiel UCU1 (strikte Konvexität), UCU2 (Existenz eines Minimierers) und UCU3 (quantitative quadratische untere Schranke) wird als technischer Kern formuliert. Die Renormierungsgruppe (RG) liefert die mathematische Brücke zwischen den Skalen. Wir tabellieren den Status von UCU1/UCU2/UCU3 über repräsentative Anwendungen hinweg (Thermodynamik, Chemie, Biologie, K41-Turbulenz, kosmologische Dunkle Energie, Yang--Mills Masselücke) und dokumentieren die aus der kosmologischen Route (\emph{Pattern~A}) gelernte methodische Lektion: Eine Positivitäts-Umformulierung einer offenen Vermutung ist im Allgemeinen eher ein Wechsel der formalen Sprache als ein Fortschritt, sofern ihre Beweisstärke nicht unabhängig gerechtfertigt ist. Dieser Artikel ist eine \emph{ehrliche Bestandsaufnahme}: Er trennt etablierte strukturelle Beiträge von bedingten Behauptungen und verzeichnet die offenen Teilprobleme (HD3' für Beal, volles UCU3 für Yang--Mills, L4 RG-Matching für die Dunkle-Energie/CRM-VI-Linie).
\end{abstract}

\textbf{Schlüsselwörter:} Stabilität, Entropieproduktion, Renormierungsgruppe,
Nash-Gleichgewicht, Skaleninvarianz, Fine-Tuning, Dunkle Energie,
Turbulenz, Masselücke.

\clearpage

% ===================================================================
\section{Einleitung}
\label{sec:intro}
% ===================================================================

Das Fine-Tuning-Problem -- die empirische Beobachtung, dass eine breite Palette von physikalischen Parametern (Standardmodell-Kopplungen, die kosmologische Konstante, die mit Leben kompatiblen chemischen Reaktionsnetzwerke) mit hoher Präzision unter allen logisch möglichen Werten ausgewählt zu sein scheint -- wird üblicherweise als skalen-spezifisch behandelt: anthropisch auf kosmologischer Skala, kinetisch auf chemischer Skala, evolutionär auf biologischer Skala. Die Functional Stability Theory (FST) untersucht die alternative Hypothese, dass diese Auswahlaussagen unterschiedliche Manifestationen eines einzigen Stabilitätsprinzips sind.

Das Programm hat über alle Skalen hinweg eine feste Form:
\begin{enumerate}[label=(\roman*),leftmargin=2em]
\item Ein Konfigurationsraum $X$, auf dem ein Funktional $F$ definiert ist.
\item Ein bedingtes System: $F$ lässt einen eindeutigen Minimierer zu, vorausgesetzt, eine strikte Konvexitätseigenschaft (UCU1), eine Existenzeigenschaft (UCU2) und eine quantitative quadratische untere Schranke (UCU3) gelten auf einer abgeschlossenen konvexen zulässigen Menge $K \subset X$.
\item Eine Renormierungsgruppen-(RG)-Brücke, die zwei Skalen verbindet und sicherstellt, dass der Minimierer auf einer Skala eine Stabilitätsaussage auf einer anderen induziert.
\end{enumerate}
Der Beitrag dieses Überblicks ist daher nicht ein neues Theorem, sondern der \emph{programmatische} Vorschlag, dass so unterschiedliche Skalen wie die Parameterselektion im Standardmodell (FST-I) und die spätzeitliche Beschleunigung durch Dunkle Energie (CRM-VI) dieselbe bedingte Struktur zulassen. Frühere Arbeiten in der Reihe behandeln einzelne Skalen; der vorliegende Artikel bildet sie auf eine einzige Architektur ab und identifiziert, welche Teilprobleme offen bleiben.

% ===================================================================
\subsection*{Die FST-Serie im Überblick}
\label{subsec:series-glance-de}
% ===================================================================

Dieser Artikel dient als \emph{programmatischer Hub} der Functional-Stability-Theory-Serie. Tabelle~\ref{tab:series-components} listet die direkten Anwendungs-Skalen-Begleiter (FST-I, FST-II, FST-III, FST-DE) sowie die wichtigsten publizierten FST-Bausteine in der Mathematik- und Physik-Linie, auf die dieser Überblick verweist.

\begin{table}[h]
\centering
\renewcommand{\arraystretch}{1.18}
\scriptsize
\begingroup
\setlength{\tabcolsep}{4pt}
\begin{tabularx}{\textwidth}{p{3.0cm}Y p{2.3cm}p{1.8cm}}
\toprule
\textbf{Beitrag} & \textbf{Skala / Ziel} & \textbf{Referenz} &
\textbf{Rolle in der Serie} \\
\midrule
\multicolumn{4}{l}{\textit{Direkte Anwendungs-Begleiter (FST-PRACTICAL)}} \\
FST-I   & Thermodynamisch / SM-Parameter      & \cite{GeigerFST-I-2026}     & Teil \\
FST-II  & Chemisch / Autokatalyse             & \cite{GeigerFST-II-2026}    & Teil \\
FST-III & Biologisch / Proteinfaltung         & \cite{GeigerFST-III-2026}   & Teil \\
FST-IV  & Kosmologischer Sammler-Slot         & in Vorbereitung             & Teil \\
\quad $\hookrightarrow$ FST-DE & Teilbeitrag: Dunkle Energie & \cite{GeigerFST-DE-2026} & Sub-Komponente \\
\midrule
\multicolumn{4}{l}{\textit{Mathematische Bausteine (FST-MATHEMATICS)}} \\
RH / Even Dominance             & spektrale Stabilität    & \cite{GeigerRH-EvenDom-2026} & Verweis \\
Zeta Zoo (FST-Mathematik)       & Klassifikation          & \cite{GeigerZetaZoo-2026}    & Verweis \\
Zookeeper (CCM-Route)           & RH technischer Kern     & \cite{GeigerZookeeper-2026}  & Verweis \\
FST Spectrum Duality            & Physics-Master          & \cite{GeigerSpectrumDuality-2026} & Verweis \\
Hodge Positivity                & AP / Algebraizität      & \cite{GeigerHodge-2026}      & Verweis \\
BSD Positivity                  & Néron--Tate-Höhen       & \cite{GeigerBSD-2026}        & Verweis \\
P vs NP Entropy                 & Komplexität             & \cite{GeigerPvsNP-2026}      & Verweis \\
\midrule
\multicolumn{4}{l}{\textit{Physikalische Bausteine (FST-PHYSICS)}} \\
Yang--Mills Masselücke          & Strong-Coupling-Gitter  & \cite{GeigerYM-2026}         & Verweis \\
Navier--Stokes Regularität      & 3D-NS-Attraktor-Route   & \cite{GeigerNS-2026}         & Verweis \\
NS-LDI-Rahmen                   & TLL/LDI-Brücke          & \cite{GeigerNSLDI-2026}      & Verweis \\
K41-Turbulenz (Joint Minimiser) & Kolmogorov-Spektrum     & \cite{GeigerK41-2026}        & Verweis \\
Turbulenz / DFC-Kaskadenroute   & anomale Dissipation     & \cite{GeigerTurbulence-2026} & Verweis \\
\bottomrule
\end{tabularx}
\endgroup
\caption{FST-Serien-Komponenten. ``Teil'' (Zenodo Related Identifier
\texttt{hasPart}) markiert die vier direkten Anwendungs-Begleiter, die
parallel zu diesem Überblick publiziert werden. ``Verweis'' (Zenodo Related Identifier
\texttt{references}) markiert publizierte FST-Bausteine, auf die dieser Überblick verweist, die aber
nicht formell mit ihm gebündelt werden. Alle Identifier sind persistente
Zenodo-DOIs; siehe die Bibliographie für die Resolver.}
\label{tab:series-components}
\end{table}

% ===================================================================
\section{Das universelle Spiel}
\label{sec:universal-game}
% ===================================================================

\subsection{Zwei Axiome}

\begin{axiom}[Anfangsbedingung]
Der Kosmos beginnt in einem Zustand von maximalem Gradienten und minimaler Struktur (niedrige Entropie, hoher Freie-Energie-Gradient).
\end{axiom}

\begin{axiom}[Entwicklung: MEPP]
Offene Nichtgleichgewichtssubsysteme entwickeln sich auf den interessierenden Zeitskalen entlang von Trajektorien, die mit dem Prinzip der maximalen Entropieproduktion (MEPP) konsistent sind~\cite{MartyushevSeleznev2006}.
\end{axiom}

\subsection{Skalen-Tripel}

Jeder physikalischen Skala $L$ ordnen wir ein Tripel zu:
\[
\mathcal{S}_L \;=\; (x_L,\; G_L,\; D_L),
\]
wobei $x_L$ die Konfigurationsvariable ist (Feld, Populationsanteil, Dichteprofil), $G_L = (S_L, \Pi_L)$ das lokale Spiel (Strategieraum $S_L$ und Auszahlung $\Pi_L$), und $D_L$ die Dynamik $\dot{x}_L = D_L(x_L; \Pi_L)$.

\subsection{Stabilitätsdefinition}

\begin{definition}[Stabiler Zustand]
$x_L^*$ ist \emph{stabil} auf der Skala $L$, wenn es ein lokal attraktiver Fixpunkt von $D_L$ (im Sinne von Lyapunov) und robust gegenüber Invasionen oder Störungen ist (Nash/ESS-Kriterium auf der Spielebene).
Diese Nash/ESS-Sprache folgt dem Standardhintergrund der evolutionären Spieltheorie nach Maynard Smith~\cite{MaynardSmith1982}.
\end{definition}

Die Wahl des Auszahlungs-Proxys auf jeder Skala (Entropieproduktion $\dot S$, Zerstörung freier Energie $-\Delta G$, mittlere Fitness $\bar f$, etc.) wird durch den physikalischen Rahmen diktiert; die FST-Behauptung ist, dass der \emph{langzeit-stabile} Attraktor mit dem maximalen Auszahlungs-Proxy unter MEPP-kompatiblen Randbedingungen korreliert.
Die kosmologische Linie nutzt außerdem die thermodynamische Raumzeit-Intuition, wie sie Jacobsons Lesart der Einstein-Gleichung als Zustandsgleichung exemplarisch formuliert~\cite{Jacobson1995}.

% ===================================================================
\section{Das universelle Stabilitätskriterium (UCU)}
\label{sec:ucu}
% ===================================================================

Der technische Kern des Programms ist ein Drei-Lemma-Endspiel, formuliert in der konvexanalytischen Sprache reflexiver Banachräume. Wir verweisen auf \texttt{CORE/UCU\_FORMAL\_DEFINITIONS.md} für die formalen Definitionen und die Standardliteratur (Brezis~\cite{Brezis2011}, Ekeland--T\'emam~\cite{EkelandTemam1999}, Evans~\cite{Evans2010}). In diesem Abschnitt ist $X$ durchgehend ein reflexiver reeller Banachraum, $K \subset X$ ist nicht-leer, abgeschlossen und konvex, und $F : K \to \mathbb{R} \cup \{+\infty\}$ erfüllt (A1) Eigentlichkeit, (A2) schwache Unterhalbstetigkeit, (A3) Koerzivität.

\begin{lemma}[UCU1 -- Strikte Konvexität]
\label{lem:ucu1}
$F$ ist \emph{strikt konvex} auf $K$:
$F(\lambda x + (1{-}\lambda) y) < \lambda F(x) + (1{-}\lambda)F(y)$
für alle $x \neq y$ in $K$ und $\lambda \in (0,1)$.
\end{lemma}

\begin{lemma}[UCU2 -- Existenz]
\label{lem:ucu2}
Unter (A1)--(A3) nimmt $F$ sein Infimum auf $K$ an:
es existiert ein $x^* \in K$ mit $F(x^*) = \inf_K F$.
\end{lemma}

\begin{lemma}[UCU3 -- Quantitative quadratische untere Schranke]
\label{lem:ucu3}
Es existiert ein $\mu > 0$, so dass für alle $x \in K$ gilt:
\[
F(x) \;\ge\; F(x^*) \;+\; \tfrac{\mu}{2}\,\|x - x^*\|^2.
\]
UCU3 spaltet sich in eine \emph{lokale} und eine \emph{globale} Version auf (siehe Abschnitt~\ref{sec:cross-scale}).
\end{lemma}

\noindent\textbf{Konsequenz (UCU-Theorem, bedingte Form).}
Wenn UCU1 + UCU2 + UCU3 auf $K$ gelten, dann lässt $F$ einen \emph{eindeutigen} Minimierer zu, und das Minimum ist in $\|\cdot\|$ \emph{quantitativ isoliert}. Die ersten beiden Lemmata sind klassisch (direkte Methode der Variationsrechnung / Tonelli); UCU3 ist der anwendungsspezifische Input, der das abstrakte Resultat mit der Physik der Skala verknüpft.

\begin{remark}
Die Rolle von UCU3 ist die gleiche wie die einer \emph{Spektrallücke} im Resolventen-Bild von \cite{GeigerZookeeper-2026}: Es wandelt eine qualitative Eindeutigkeitsaussage in eine quantitative Stabilitätsabschätzung um. Skalenübergreifend besagt UCU3: ``Keine flache Richtung am Minimum''.
\end{remark}

% ===================================================================
\section{Die Renormierungsgruppe als mathematische Brücke}
\label{sec:rg-bridge}
% ===================================================================

Die RG liefert den formalen Apparat, um das universelle Spiel auf der Skala $L$ mit dem auf der Skala $L+1$ in Beziehung zu setzen, in der Wilson--Kogut-Tradition von Skalentransformationen~\cite{Wilson1974}. Es wird gefordert, dass eine Kadanoff-Blocktransformation $R_b : (x_L, G_L, D_L) \mapsto (x_{L+1}, G_{L+1}, D_{L+1})$ mindestens eine der folgenden Invarianten erhält:
\begin{enumerate}[label=(\Alph*),leftmargin=2em]
\item die Fixpunktstruktur (Anzahl und Art der stabilen Attraktoren),
\item die Geometrie der besten Antwort (Vorzeichenmuster der Jacobi-Matrix um $x_L^*$),
\item die Existenz einer Lyapunov- / Potentialfunktion.
\end{enumerate}
Wann immer (A)--(C) erhalten bleiben, induziert ein UCU-stabiler Skala-$L$-Minimierer einen UCU-stabiler Skala-$(L{+}1)$-Minimierer.

\paragraph{Kompakte Aussage.}
Die RG-Brücke ist der Mechanismus, durch den das bedingte UCU-Theorem \emph{übertragbar} wird. In der Praxis wurde die Brücke bisher nur in Fragmenten ausgearbeitet (siehe die offenen Probleme in Abschnitt~\ref{sec:open}); für das kosmologische RG-Matching des $R + \gamma R^2$ Sektors (CRM-VI) bleibt der L4-Schritt offen.

% ===================================================================
\section{Skalenübergreifender Status}
\label{sec:cross-scale}
% ===================================================================

Tabelle~\ref{tab:status} fasst den gegenwärtigen Status von UCU1, UCU2, UCU3 (lokal und global) über die sechs Anwendungsskalen des Programms hinweg zusammen.

\begin{table}[h]
\centering
\renewcommand{\arraystretch}{1.18}
\scriptsize
\begingroup
\setlength{\tabcolsep}{4pt}
\begin{tabularx}{\textwidth}{p{2.5cm}p{2.6cm}c c p{2.1cm}Y}
\toprule
\textbf{Skala} & \textbf{Funktional} & \textbf{UCU1} & \textbf{UCU2}
& \textbf{UCU3 lokal} & \textbf{UCU3 global} \\
\midrule
FST-I (Thermo / SM)        & $\Sigma[\alpha,m_e/m_p,\dots]$ &
\checkmark & \checkmark & Plateau & offen \\
FST-II (Chemie)            & $G[\text{Hyperzyklus}]$ &
\checkmark & \checkmark & lokal OK & offen \\
FST-III (Biologie / Falt.) & $F_{\text{fold}}[\phi]$ &
\checkmark & \checkmark & lokal OK & offen \\
K41 (Turbulenz)            & $F[E(k)]$ &
\checkmark & \checkmark & \checkmark $(g^*=\tfrac12)$ & LSI offen \\
DE / CRM-VI (Kosmologie)   & $\Phi(\rho)$ &
\checkmark & \checkmark & \checkmark & L4 RG-Match offen \\
YM (Masselücke)            & $\mathcal{E}[A]$ &
? & ? & ? & UCU3 = \emph{Masselücken-Frage} \\
\bottomrule
\end{tabularx}
\endgroup
\caption{Status von UCU1/UCU2/UCU3 über verschiedene Anwendungsskalen.
\checkmark{} = bedingt etabliert unter Standardvoraussetzungen (Brezis,
Ekeland--T\'emam direkte Methode); ``offen'' = Teilproblem identifiziert,
aber noch nicht innerhalb des FST-Rahmens gelöst. Die Yang--Mills Linie
ist beigefügt, um die Analogie zu markieren: UCU3 (global) für das YM-Funktional
\emph{ist} die Masselücken-Frage selbst.}
\label{tab:status}
\end{table}

\paragraph{Lesehilfe zur Tabelle.}
Für Thermodynamik, Chemie und Biologie ergeben sich UCU1+UCU2 aus der Konvexität von Funktionalen (ähnlich Freier Energie) auf physikalisch motivierten zulässigen Mengen; UCU3 ist lokal, aber nicht global, weil Nichtgleichgewichtsrandbedingungen konkurrierende Anziehungsbereiche (Basins) erzeugen können. Für die K41-Turbulenz wird das lokale UCU3 beim Kolmogorov-Skalierungsexponenten $g^*=\tfrac12$ erreicht, aber eine logarithmische Sobolev-Ungleichung (LSI) wird benötigt, um zu einer globalen Aussage hochzustufen. Für Dunkle Energie (CRM-VI) gelten UCU1--UCU3 für das de~Sitter-Funktional $\Phi(\rho)$ auf $\mathbb{R}_+$, aber die Hochstufung des Resultats durch das L4 RG-Matching zu einer vollständigen effektiven Feldtheorie ist noch nicht abgeschlossen.

\subsection{Status der FST-Bausteine (Mathematik und Physik)}
\label{subsec:building-blocks}

Über die UCU-Anwendungstabelle hinaus umfasst das FST-Programm eine Reihe
publizierter \emph{Bausteine} in der Mathematik- und Physik-Linie.
Tabelle~\ref{tab:building-blocks} fasst deren ehrlichen Beweisstatus
zusammen (entnommen den individuellen Beweisnotizen der jeweiligen
Projekte, siehe \texttt{STATUS\_UEBERSICHT.md} im Repository).

\begin{table}[h]
\centering
\renewcommand{\arraystretch}{1.18}
\scriptsize
\begingroup
\setlength{\tabcolsep}{4pt}
\begin{tabularx}{\textwidth}{p{2.6cm}p{1.8cm}p{2.4cm}Y}
\toprule
\textbf{Baustein} & \textbf{Linie} & \textbf{Status} &
\textbf{Offene Kernlücke} \\
\midrule
RH / Even Dominance      & Mathematik & unbedingt   & analytische
$c_{\text{gap}}$-Ableitung entfernt letzten numerischen Input \\
Zeta Zoo / FST-Mathematik & Mathematik & bedingt     & Prime-Hub /
OP5-Obstruktionen \\
Hodge                    & Mathematik & Framework   & Hard Direction
(AP $\Rightarrow$ algebraisch) \\
BSD                      & Mathematik & Framework   & Higher
Gross--Zagier für Rang $\ge 2$ \\
P vs NP (Entropie)       & Mathematik & Reformulierung & Uniformity
Bridge (KEIN Beweis) \\
\midrule
Yang--Mills Masselücke   & Physik     & Strong-Coupling bewiesen &
RG/Continuum-Transport bedingt \\
Navier--Stokes (Attraktor)& Physik    & bedingtes Framework &
Attraktor-Regularitätspaket, Assumption G/G' \\
NS-LDI                   & Physik     & Framework + Proof of Life &
TLL+LDI für 3D-NS-Attraktoren \\
K41-Turbulenz (Paper A)  & Physik     & unbedingt   & --- (Paper B
bleibt bedingt auf $\mathrm{DFC1}^{\vee}$) \\
Turbulenz / DFC-Kaskade  & Physik     & bedingtes Framework &
anomale Dissipation hängt von $\mathrm{DFC1}^{\vee}$ und der DFC-Brückenschließung ab \\
\midrule
FST-DE (Dunkle Energie)  & Kosmologie & Framework + No-Go &
UV-Renormierung, Screening-Lift \\
\bottomrule
\end{tabularx}
\endgroup
\caption{Ehrlicher Beweisstatus der FST-Bausteine in der Mathematik-,
Physik- und Kosmologie-Linie (Stand 2026-05). Jeder Eintrag verweist auf
ein eigenes Projekt, dessen primäre Beweisnotiz (BEWEISNOTIZ.md) die
genaue Lemma-Kette und die verbliebenen offenen Lücken dokumentiert.
Der vorliegende Überblick beweist oder erweitert keines dieser Resultate;
er verortet jedes innerhalb derselben UCU/RG/MEPP-Architektur. Wo der
Status ``Reformulierung'' (P vs NP) oder ``Framework'' (Hodge, BSD,
NS, FST-DE) lautet, ist die offene Kernlücke explizit die ursprüngliche
Millennium-artige Schwierigkeit -- konsistent mit der Pattern-A-Lektion
(Abschnitt~\ref{sec:pattern-a}).}
\label{tab:building-blocks}
\end{table}

% ===================================================================
\section{Die Pattern-A-Lektion}
\label{sec:pattern-a}
% ===================================================================

Sowohl in der mathematischen Linie (Hodge, Beal/abc, BSD, P~vs~NP) als auch in der kosmologischen Linie (CRM-VI, Dunkle Energie) tauchte ein methodisches Muster so oft auf, dass wir es als programmweite \emph{Warnung} behandeln.

\begin{quote}
\textbf{Pattern A (negative Form).}
Eine Umformulierung einer offenen Vermutung $V$ in eine ``Positivitäts-Normalform'' $P$ ist im Allgemeinen nur eine Übersetzung der Sprache: Typischerweise gilt $P \Leftrightarrow V$. Ohne einen unabhängigen Stärkenvergleich ist $P$ kein Beweis für $V$.
\end{quote}

\noindent Die kosmologische Variante ist operativ identisch: Die Abbildung einer \emph{physikalischen} Frage $V_{\text{phys}}$ auf ein Modell $M$, das eine offene mathematische Teilfrage $S$ enthält, löst $V_{\text{phys}}$ nicht; es reduziert sie lediglich bedingt auf $S$. Das kosmologische CC-Problem (Problem der kosmologischen Konstanten) und das inflationäre Anfangswertproblem trafen in der CRM-VI-Linie jeweils auf dieses Muster. Die Lektion ist in \texttt{CORE/PATTERN\_A\_LESSON.md} und der begleitenden internen Notiz~\cite{FST-Unified2026} festgehalten; sie ist der Grund, warum jeder UCU-Anwendungsbeitrag nun einen Abschnitt namens \emph{Ehrlicher Status} enthält.

% ===================================================================
\section{Offene Teilprobleme}
\label{sec:open}
% ===================================================================

Das Programm hinterlässt drei scharf umrissene, offene Probleme.

\paragraph{HD3' für Beal.}
Der Beal-Verifier-Sieve-Rahmen reduziert die Beal-Vermutung auf vier Teilaussagen HD3'-A bis HD3'-D. Wie in \texttt{HD3\_B\_EFFECTIVE\_CHARACTER\_SUMS.md} und \texttt{HD3\_ACD\_TRIO\_OBSTRUCTIONS.md} dokumentiert, liegt die starke Charakter-Summen-Schranke HD3'-B (Kürzung $q^{1+\varepsilon}$ statt der üblichen Wurzel-Kürzung $q^{3/2}$) jenseits der üblichen Weil--Deligne-Werkzeuge, und HD3'-A (Salt-Closing) sowie HD3'-C (Counting) lassen sich darauf reduzieren. Eine Notiz, die die bedingte Struktur von HD3'-A unter Annahme von HD3'-B dokumentiert, ist in \texttt{HD3\_A\_SALT\_CLOSING.md} verfügbar.

\paragraph{Globales UCU3 für Yang--Mills.}
Die Masselücken-Frage für reine $4d$ Yang--Mills-Theorie ist genau die Frage, ob das Energiefunktional $\mathcal{E}[A]$ eine \emph{globale} quadratische untere Schranke um das Vakuum zulässt: ein globales UCU3. Diese Identifizierung bedeutet, dass der FST-Rahmen Yang--Mills nicht löst; er verortet die Masselücken-Frage lediglich als den kanonischen UCU3-globalen Repräsentanten. UCU1 und UCU2 im YM-Kontext sind innerhalb der FST-Sprache selbst offene Fragen.

\paragraph{L4 RG-Matching für CRM-VI.}
Die kosmologische Dunkle-Energie-Linie (CRM-VI), entwickelt im CRM-Programm~\cite{GeigerCRM2026}, liefert UCU1--UCU3 auf einem reduzierten Konfigurationsraum (de~Sitter-Funktional in $\mathbb{R}_+$). Das Anheben des Resultats durch ein vierstufiges RG-Matching (L1 mikroskopische Wirkung, L2 effektives Potential, L3 kosmologischer Hintergrund, L4 Störungen + Abgleich mit Daten) lässt L4 noch offen. Dies ist das kosmologische Analogon zur Pattern-A-Lektion: Die mathematische Reduktion ist bedingt durch ein physikalisches RG-Matching, das noch nicht auf dem für eine Publikation erforderlichen Niveau durchgeführt wurde.

% ===================================================================
\section{Ehrlicher Status}
\label{sec:honest-status}
% ===================================================================

\begin{description}[leftmargin=2em,style=nextline]
\item[Was FST liefert.]
Eine gemeinsame Architektur (universelles Spiel $+$ MEPP-Axiom $+$ UCU Drei-Lemma-Endspiel $+$ RG-Brücke), unter der Skalen wie der Parameterraum des Standardmodells und die spätzeitliche Beschleunigung durch Dunkle Energie eine einzige bedingte Eindeutigkeits-/Stabilitätsaussage zulassen. Eine Statuskarte (Tabelle~\ref{tab:status}), die zeigt, wo die Architektur geschlossen und wo sie bedingt ist. Eine methodische Lektion (Pattern~A), die bereits mehrfache Behauptungen von falschem Fortschritt in der mathematischen Linie verhindert hat.

\item[Was FST \emph{nicht} liefert.]
Eine Lösung eines individuellen harten, offenen Problems (Yang--Mills Masselücke, Beal, BSD). FST legt die \emph{Form} der bedingten Aussage fest; der anwendungsspezifische globale UCU3-Schritt wird in jedem harten Fall mit der ursprünglichen Schwierigkeit identifiziert. Die Neuerung ist daher architektonisch und reduktiv, nicht konstruktiv.

\item[Falsifikationskriterium.]
Das Programm ist falsifiziert, wenn für eine der aufgelisteten Skalen UCU1 oder UCU2 auf der natürlichen zulässigen Menge scheitert. Ein Scheitern von UCU3-global ist keine Falsifikation: Es markiert das offene Teilproblem auf dieser Skala.
\end{description}

% ===================================================================
\section*{Rechenhinweis}
% ===================================================================
Dieser Artikel ist ein Programmüberblick. Berechnungen zur Unterstützung der kosmologischen Linie (Watkins-artige Spektrum-Experimente) wurden auf \texttt{ellmos-services} durchgeführt; der vorliegende Überblick selbst führt kein neues numerisches Experiment ein.

% ===================================================================
\section*{KI-Offenlegung (Stufe L3)}
% ===================================================================
Dieser Beitrag wurde unter substantieller KI-Unterstützung (Anthropic Claude, OpenAI GPT, Google Gemini in verschiedenen Varianten) erstellt. KI-Werkzeuge wurden eingesetzt für: Synthese über mehrere Beiträge hinweg zwischen den FST-I/II/III/DE-Begleitern; strukturelle Bearbeitung und Konsistenzprüfung der UCU/RG-Architektur; Generierung und Verifikation der skalenübergreifenden Statuskarten (Tabellen~\ref{tab:series-components}, \ref{tab:status} und~\ref{tab:building-blocks}); Literatur-Querverweise mit Zenodo-DOIs; sowie bilinguale EN/DE-Synchronisation. Alle wissenschaftlichen Aussagen, die UCU-Drei-Lemma-Architektur, die Pattern-A-Lektion und das Falsifikationskriterium sind die eigenen Aussagen des Autors und unterliegen der üblichen Autorenverantwortung. Kein Abschnitt wurde autonom von einer KI geschrieben. Die Kategorisierung folgt dem fünfstufigen KI-Offenlegungsschema des FST-Programms (L1 = keine KI; L5 = KI als hauptverantwortlich); dieser Beitrag ist L3 (KI als substantielle Assistenz, Mensch als hauptverantwortlicher Autor).

% ===================================================================
% Bibliographie (Skeleton -- wird in v1 vervollständigt)
% ===================================================================
\begin{thebibliography}{99}

\bibitem{MartyushevSeleznev2006}
L.M.~Martyushev, V.D.~Seleznev,
\emph{Maximum entropy production principle in physics, chemistry and biology},
Phys.~Rep.~\textbf{426} (2006), 1--45.
DOI: \url{https://doi.org/10.1016/j.physrep.2005.12.001}.

\bibitem{Brezis2011}
H.~Brezis,
\emph{Functional Analysis, Sobolev Spaces and Partial Differential Equations},
Springer, 2011.
DOI: \url{https://doi.org/10.1007/978-0-387-70914-7}.

\bibitem{EkelandTemam1999}
I.~Ekeland, R.~T\'emam,
\emph{Convex Analysis and Variational Problems},
SIAM Classics, 1999.
DOI: \url{https://doi.org/10.1137/1.9781611971088}.

\bibitem{Evans2010}
L.C.~Evans,
\emph{Partial Differential Equations}, 2nd~ed.,
Amer.~Math.~Soc., 2010.
DOI: \url{https://doi.org/10.1090/gsm/019}.

\bibitem{Wilson1974}
K.G.~Wilson, J.~Kogut,
\emph{The renormalization group and the $\varepsilon$ expansion},
Phys.~Rep.~\textbf{12} (1974), 75--199.
DOI: \url{https://doi.org/10.1016/0370-1573(74)90023-4}.

\bibitem{Jacobson1995}
T.~Jacobson,
\emph{Thermodynamics of spacetime: the Einstein equation of state},
Phys.~Rev.~Lett.~\textbf{75} (1995), 1260--1263.
DOI: \url{https://doi.org/10.1103/PhysRevLett.75.1260}.

\bibitem{MaynardSmith1982}
J.~Maynard~Smith,
\emph{Evolution and the Theory of Games},
Cambridge Univ.~Press, 1982.
DOI: \url{https://doi.org/10.1017/CBO9780511806292}.

% --- FST-Programm-Begleiter (hasPart) -------------------------------
\bibitem{GeigerFST-I-2026}
L.~Geiger,
\emph{Functional Stability Theory I: A Game-Theoretic Framework
for the Thermodynamic Stability of Fundamental Parameters},
Zenodo: \texttt{10.5281/zenodo.20130544}, 2026.

\bibitem{GeigerFST-II-2026}
L.~Geiger,
\emph{Functional Stability Theory II: Chemical Stability and
Autocatalytic Selection},
Zenodo: \texttt{10.5281/zenodo.20130563}, 2026.

\bibitem{GeigerFST-III-2026}
L.~Geiger,
\emph{Functional Stability Theory III: Biological Stability and
Nash Frustration},
Zenodo: \texttt{10.5281/zenodo.20130573}, 2026.

\bibitem{GeigerFST-DE-2026}
L.~Geiger,
\emph{Dark Energy as Residual Vacuum Free Energy: A Thermodynamic
Bound on the Cosmological Constant},
Zenodo: \texttt{10.5281/zenodo.19036235}, 2026.

% --- Mathematische FST-Bausteine (references) ----------------------
\bibitem{GeigerRH-EvenDom-2026}
L.~Geiger,
\emph{From Landscape to Proof: The Even-Dominance Route to the
Riemann Hypothesis},
Zenodo: \texttt{10.5281/zenodo.19035640}, 2026.

\bibitem{GeigerZetaZoo-2026}
L.~Geiger,
\emph{The Zeta Zoo: The Mathematical Side of Functional Stability Theory},
Zenodo: \texttt{10.5281/zenodo.19673226}, 2026.

\bibitem{GeigerZookeeper-2026}
L.~Geiger,
\emph{The Spectral Zookeeper: The Riemann Hypothesis via Spectral
Microcluster Closure in the Connes-Consani-Moscovici Framework},
Zenodo: \texttt{10.5281/zenodo.19673126}, 2026.

\bibitem{GeigerSpectrumDuality-2026}
L.~Geiger,
\emph{FST Spectrum Duality: The Renormalized Free Energy Principle as
Physical Instantiation of Functional Stability Theory},
Zenodo: \texttt{10.5281/zenodo.19036190}, 2026.

\bibitem{GeigerHodge-2026}
L.~Geiger,
\emph{Arithmetic Positivity and Absolute Hodge Classes: Why HR
Positivity Does Not Strengthen Deligne's Absoluteness},
Zenodo: \texttt{10.5281/zenodo.19087439}, 2026.

\bibitem{GeigerBSD-2026}
L.~Geiger,
\emph{The BSD Conjecture as a Positivity Normal-Form Theorem},
Zenodo: \texttt{10.5281/zenodo.19087443}, 2026.

\bibitem{GeigerPvsNP-2026}
L.~Geiger,
\emph{An Entropic Perspective on P vs NP: Reformulation via Kolmogorov
Complexity},
Zenodo: \texttt{10.5281/zenodo.19056809}, 2026.

% --- Physikalische FST-Bausteine (references) ----------------------
\bibitem{GeigerYM-2026}
L.~Geiger,
\emph{Volume-Independent Spectral Gap for Strong-Coupling Lattice
Gauge Theory via Poincar\'e-Dobrushin Machinery},
Zenodo: \texttt{10.5281/zenodo.19087433}, 2026.

\bibitem{GeigerNS-2026}
L.~Geiger,
\emph{A Thermodynamic Obstruction to Finite-Time Blow-Up in the
3D Navier--Stokes Equations},
Zenodo: \texttt{10.5281/zenodo.19087449} (Concept), 2026.

\bibitem{GeigerNSLDI-2026}
L.~Geiger,
\emph{Log-Distance Integrability for Dissipative Flows: BV Selections
on Fractal Attractors},
Zenodo: \texttt{10.5281/zenodo.19056807}, 2026.

\bibitem{GeigerK41-2026}
L.~Geiger,
\emph{The K41 Spectrum as the Unique Variational Minimiser of a
Scale-Resolved Free Energy},
Zenodo: \texttt{10.5281/zenodo.20131305}, 2026.

\bibitem{GeigerTurbulence-2026}
L.~Geiger,
\emph{Anomalous Dissipation and the Turbulent Cascade: A Variational
Free-Energy Characterisation of Kolmogorov Scaling},
Zenodo: \texttt{10.5281/zenodo.19056813}, 2026.

% --- Ältere Serien und interne Notizen -----------------------------
\bibitem{GeigerCRM2026}
L.~Geiger,
\emph{The Curvature Relaxation Model: A Four-Paper Program for
Geometric Cosmology Without the Dark Sector},
Zenodo: \texttt{10.5281/zenodo.18728935}, 2026.

\bibitem{FST-Unified2026}
L.~Geiger,
\emph{Pattern A as the Recurring Stability Logic across the FST
Research Ecosystem}, interne Notiz (CORE/PATTERN\_A\_LESSON.md),
2026.

\end{thebibliography}

\end{document}
